\textbf{Phần 1. Trắc nghiệm lựa chọn}
\setcounterpageref{ex}{0}

\begin{ex} %[2D2V10]
Tính giá trị: $7 + 9 - 8$.
\choice
{$7$}
{$9$}
{$6$}
{\True $8$}
\loigiai{
$7 + 9 = 16$; $16 - 8 = 8$.
\textbf{Đáp án của câu hỏi là D}
}
\end{ex}

\begin{ex} %[2D2V10]
Một lớp học có $9$ bạn nam và $8$ bạn nữ. Sau đó có thêm $2$ bạn nữ chuyển vào lớp. Hỏi lớp học lúc này có bao nhiêu học sinh?
\choice
{$18$ học sinh}
{\True $19$ học sinh}
{$20$ học sinh}
{$17$ học sinh}
\loigiai{
Tổng số học sinh lúc này: $9 + 8 + 2 = 9 + 10 = 19$ (học sinh).
\textbf{Đáp án của câu hỏi là B}
}
\end{ex}

\begin{ex} %[2D2V10]
Tìm $x$ biết: $x + 6 = 8 + 7$.
\choice
{$8$}
{$10$}
{\True $9$}
{$7$}
\loigiai{
Vế phải: $8 + 7 = 15$. Ta có: $x = 15 - 6 = 9$.
\textbf{Đáp án của câu hỏi là C}
}
\end{ex}

\begin{ex} %[2D2V10]
Bạn Mai có $7$ cái kẹo, bạn Hoa có nhiều hơn Mai $4$ cái kẹo. Hỏi cả hai bạn có tất cả bao nhiêu cái kẹo?
\choice
{\True $18$ cái kẹo}
{$17$ cái kẹo}
{$11$ cái kẹo}
{$19$ cái kẹo}
\loigiai{
Hoa có: $7 + 4 = 11$ (cái kẹo).\nCả hai bạn có: $7 + 11 = 18$ (cái kẹo).
\textbf{Đáp án của câu hỏi là A}
}
\end{ex}

\begin{ex} %[2D2V10]
Điền dấu thích hợp: $8 + 8 \dots 9 + 7$.
\choice
{$>$}
{$<$}
{Không so sánh được}
{\True $=$}
\loigiai{
$8 + 8 = 16$; $9 + 7 = 16$. Hai vế bằng nhau nên điền dấu $=$.
\textbf{Đáp án của câu hỏi là D}
}
\end{ex}

\begin{ex} %[2D2V10]
Tìm một số biết rằng lấy số đó cộng với $5$ rồi trừ đi $8$ thì được $6$.
\choice
{$8$}
{\True $9$}
{$10$}
{$7$}
\loigiai{
Tính ngược từ cuối: $6 + 8 - 5 = 14 - 5 = 9$.
\textbf{Đáp án của câu hỏi là B}
}
\end{ex}

\begin{ex} %[2D2V10]
Có hai bao gạo, bao thứ nhất nặng $9\text{ kg}$, bao thứ hai nặng $8\text{ kg}$. Người ta đã ăn hết $5\text{ kg}$. Hỏi còn lại bao nhiêu ki-lô-gam gạo?
\choice
{$11\text{ kg}$}
{$13\text{ kg}$}
{\True $12\text{ kg}$}
{$10\text{ kg}$}
\loigiai{
Tổng số gạo: $9 + 8 = 17\text{ (kg)}$. Còn lại: $17 - 5 = 12\text{ (kg)}$.
\textbf{Đáp án của câu hỏi là C}
}
\end{ex}

\begin{ex} %[2D2V10]
Cho ba số: $6, 8, 9$. Tổng của hai số lớn nhất trong ba số trên là
\choice
{\True $17$}
{$15$}
{$14$}
{$18$}
\loigiai{
Hai số lớn nhất là 8 và 9. Tổng là: $8 + 9 = 17$.
\textbf{Đáp án của câu hỏi là A}
}
\end{ex}

\begin{ex} %[2D2V10]
Một đàn vịt đang bơi dưới ao có $9$ con, trên bờ có $6$ con. Sau đó có $3$ con từ dưới ao nhảy lên bờ. Hỏi lúc này dưới ao còn bao nhiêu con vịt?
\choice
{$7$ con vịt}
{$5$ con vịt}
{$8$ con vịt}
{\True $6$ con vịt}
\loigiai{
Số vịt dưới ao lúc này: $9 - 3 = 6$ (con vịt).
\textbf{Đáp án của câu hỏi là D}
}
\end{ex}

\begin{ex} %[2D2V10]
Tổng của ba số bằng $19$. Biết hai số đầu là $7$ và $8$. Số thứ ba là
\choice
{$3$}
{\True $4$}
{$5$}
{$6$}
\loigiai{
Tổng hai số đầu là: $7 + 8 = 15$. Số thứ ba là: $19 - 15 = 4$.
\textbf{Đáp án của câu hỏi là B}
}
\end{ex}

\begin{ex} %[2D2V10]
Hai số tự nhiên có tổng bằng $13$ và hiệu bằng $3$. Hai số đó là
\choice
{$9\text{ và } 4$}
{$7\text{ và } 6$}
{\True $8\text{ và } 5$}
{$10\text{ và } 3$}
\loigiai{
Số lớn là: $(13 + 3) : 2 = 8$. Số bé là: $8 - 3 = 5$.
\textbf{Đáp án của câu hỏi là C}
}
\end{ex}

\begin{ex} %[2D2V10]
Điền số thích hợp: $9 + 5 = \dots + 8$.
\choice
{\True $6$}
{$5$}
{$7$}
{$8$}
\loigiai{
Vế trái: $9 + 5 = 14$. Ta có: $14 - 8 = 6$. Vậy điền 6.
\textbf{Đáp án của câu hỏi là A}
}
\end{ex}

\textbf{Phần 2. Trắc nghiệm đúng sai}
\setcounterpageref{ex}{0}

\begin{ex} %[2D2V10]
Xét tính đúng sai của các bài toán tổng hợp cộng trong phạm vi 20:
\choiceTF
{\True Biểu thức $8 + 7 + 3$ có kết quả bằng $18$.}
{Biểu thức $9 + 6$ có kết quả bé hơn $8 + 6$.}
{\True Nếu thêm 2 đơn vị vào một số hạng thì tổng tăng thêm 2 đơn vị.}
{Tổng của hai số có 1 chữ số có thể bằng 20.}
\loigiai{
\begin{itemize}
    \item Ý a) Đúng: $8 + (7 + 3) = 8 + 10 = 18$.
    \item Ý b) Sai: $9 + 6 = 15 > 8 + 6 = 14$.
    \item Ý c) Đúng: Đúng theo tính chất phép cộng.
    \item Ý d) Sai: Tổng lớn nhất là $9 + 9 = 18 < 20$.
\end{itemize}
\begin{itemize}
    \item a) \textbf{Đúng}.
    \item b) \textbf{Sai}.
    \item c) \textbf{Đúng}.
    \item d) \textbf{Sai}.
\end{itemize}
}
\end{ex}

\begin{ex} %[2D2V10]
Bạn Hùng có 8 viên bi đỏ và 9 viên bi xanh:
\choiceTF
{\True Hùng có tất cả 17 viên bi.}
{\True Số bi xanh nhiều hơn số bi đỏ là 1 viên bi.}
{Hùng cho bạn 8 viên bi đỏ thì Hùng còn lại 8 viên bi xanh.}
{\True Nếu Hùng nhặt thêm 3 viên bi vàng thì Hùng có tất cả 20 viên bi.}
\loigiai{
\begin{itemize}
    \item Ý a) Đúng: $8 + 9 = 17$ viên bi.
    \item Ý b) Đúng: $9 - 8 = 1$ viên bi.
    \item Ý c) Sai: Hùng có 9 viên bi xanh nên Hùng còn lại 9 viên bi xanh.
    \item Ý d) Đúng: $17 + 3 = 20$ viên bi. Đúng.
\end{itemize}
\begin{itemize}
    \item a) \textbf{Đúng}.
    \item b) \textbf{Đúng}.
    \item c) \textbf{Sai}.
    \item d) \textbf{Đúng}.
\end{itemize}
}
\end{ex}

\textbf{Phần 3. Trả lời ngắn}
\setcounterpageref{ex}{0}

\begin{ex} %[2D2V10]
Tính giá trị biểu thức: $6 + 7 + 4$.
\shortans{17}
\loigiai{
\textbf{Đáp số:} $17$\\
$(6 + 4) + 7 = 10 + 7 = 17$.
}
\end{ex}

\begin{ex} %[2D2V10]
Lan có 8 cái nhãn vở, mẹ cho Lan thêm 9 cái nhãn vở nữa. Lan tặng bạn 6 cái. Hỏi Lan còn lại bao nhiêu cái nhãn vở?
\shortans{11}
\loigiai{
\textbf{Đáp số:} $11$\\
Lan có: $8 + 9 = 17$. Còn lại: $17 - 6 = 11$ cái.
}
\end{ex}

\begin{ex} %[2D2V10]
Tìm số tự nhiên y biết: $y - 7 = 8 + 1$.
\shortans{16}
\loigiai{
\textbf{Đáp số:} $16$\\
$y - 7 = 9 \Rightarrow y = 9 + 7 = 16$.
}
\end{ex}

\begin{ex} %[2D2V10]
Trong rổ có 18 quả trứng gà và vịt. Biết số trứng gà là 9 quả. Hỏi số trứng vịt là bao nhiêu quả?
\shortans{9}
\loigiai{
\textbf{Đáp số:} $9$\\
$18 - 9 = 9$ quả.
}
\end{ex}

\begin{ex} %[2D2V10]
Tìm tổng của số lớn nhất có một chữ số và số lẻ bé nhất có một chữ số khác 1.
\shortans{12}
\loigiai{
\textbf{Đáp số:} $12$\\
Số lớn nhất có 1 chữ số là 9. Số lẻ bé nhất khác 1 là 3. Tổng: $9 + 3 = 12$.
}
\end{ex}

\begin{ex} %[2D2V10]
Tìm một số biết rằng lấy số đó cộng với 8 thì được số tròn chục nhỏ nhất có hai chữ số.
\shortans{2}
\loigiai{
\textbf{Đáp số:} $2$\\
Số tròn chục nhỏ nhất là 10. Số đó là: $10 - 8 = 2$.
}
\end{ex}
